% ============================================================
% BME VIK TDK abstract — FINAL
%
% Supersedes v1-v5 in docs/framing/. Keeps v5's opening (the
% contribution-is-a-departure observation), and rebuilds the body as
% an argument that unfolds rather than a set of parallel stages:
%
%   1. thesis + the task           (opening from v5, task line from v1)
%   2. the model that generates the abstract
%   3. WHAT IT CANNOT RECOVER — the specific methodology — and why
%      that is a property of the task, not of the model   <- the hinge
%   4. therefore the residual is the novelty
%   5. how to measure it credibly, and validate it
%
% Paragraph 3 is load-bearing. Everything before it earns it and
% everything after it depends on it. If a reviewer attacks one
% paragraph it will be that one, so the number that goes in it must be
% the best-supported number in the paper.
%
% LENGTH: ~470 words of body, cut from v5's 604 to fit one page.
%   For reference: v1 = 364 words, v2 = 464, v3 = 543, v4 = 635, v5 = 604.
% What was sacrificed, deliberately: the attribution controls (the
% text-matched baseline, the disabled and shuffled prefixes) lose
% their paragraph and survive as one clause at the end of paragraph 2.
% The full control suite belongs in the paper, not here. Also cut: the
% LLM-as-reviewer line, and the per-axis breakdown of the residual.
%
% NUMBERS: two slots. Paragraph 3 takes the recoverable-vs-
% unrecoverable figures (domain/problem/task against methodology, and
% the share of downstream node-classification performance that
% survives substitution); paragraph 5 takes the validation figures.
% Do not submit without them — the official sample requires key
% results in the body.
%
% Structure follows the official one-page sample:
% https://portal.vik.bme.hu/files/00000047.pdf
%   - max 1 page, min ~20 lines of body text
%   - NO header, footer, or page numbers
%   - figures/equations only if strictly necessary
%   - own contribution must be clearly separated from cited work
%   - body must cover: methodology, key results, conclusions
%
% NOTE ON LANGUAGE: the official TDK conference language is Hungarian.
% This English draft is for internal review with Bálint/Dani.
% Translate before submission unless an English exception is agreed
% with the section chair; the Hungarian preamble (babel[magyar], the
% author/consultant block) is already set up in abstract.tex.
% Hungarian runs longer than English — recheck the page fit after
% translating, and cut from paragraph 4's positioning clauses first.
% ============================================================

\documentclass[12pt]{article}

\usepackage[utf8]{inputenc}
\usepackage[T1]{fontenc}
\usepackage[english]{babel}
\usepackage[a4paper, margin=2.5cm]{geometry}
\usepackage{mathptmx}      % Times-like font
\usepackage{hyperref}
\usepackage{enumitem}

\pagestyle{empty}          % <- enforces "no header/footer/page numbers"
\setlength{\parindent}{0pt}
\setlength{\parskip}{0.6em}

\begin{document}

% ---------- Title ----------
\begin{center}
    {\Large\bfseries What Its Own References Cannot Anticipate: \\
    Measuring Scientific Novelty on a Citation Graph}
\end{center}

\vspace{0.8em}

% ---------- Authors ----------
\begin{center}
    Gema Csanád IV. Inf, gemacsanad@gmail.com \\

\end{center}

\vspace{0.4em}

% ---------- Consultant ----------
\begin{center}
    Consultant: Unyi Dániel
\end{center}

\vspace{1.2em}

% ==========================================================
% BODY
% ==========================================================

A paper's contribution is, by definition, its departure from the work
it cites: anything already contained in its references would not be a
contribution. We turn that observation into a measurement of
scientific novelty. In a citation graph — each node a publication,
each edge a citation from one paper to another, every node carrying
that paper's abstract as its text content — we study one task: given
a paper's reference list, but withholding the paper's own abstract
entirely, generate what that abstract says.

A graph neural network encodes the citation neighbourhood — the cited
papers, and their references in turn, beyond what could fit in a
prompt — into a continuous representation, which a trained adapter
projects into soft prefix tokens injected into the attention layers
of a language model, alongside the plain text of the directly cited
papers. The target's own features are masked and edges run from cited
to citing paper, so nothing published after it can reach it. Whether
the graph carries anything the neighbours' text does not already
convey is tested against a text-matched baseline given the same
content within the same token budget.

What this recovers is not uniform, and the pattern is the result
everything else rests on. A paper's research domain, its problem and
its task are largely determined by its bibliography and are recovered
well. Its specific methodology is not, and sits near the floor even
for far stronger models. That is not a shortcoming of the model but a
property of the task: a paper's method is precisely the part its
references do not contain. The recoverable half we measure without a
language model as judge: generated text is substituted for the real
node attributes, and we ask how much of the downstream
node-classification performance survives.

The unrecoverable half is the measurement. We take the discrepancy
between the predicted and the actual abstract as an estimate of how
much of a paper could not have been anticipated from its own
antecedents. Existing novelty measures are not conditioned on those
antecedents: reference-based scores treat a paper as a combination of
cited journals and discard its content, while perplexity-based scores
ask whether a text is unusual in general rather than given what this
paper was built on. Ours is conditional and generative — the
expectation is constructed from the paper's own reference list.

Such a residual confounds novelty with model error and with
incomplete bibliographic coverage, so it is read relative to matched
papers of the same field, year and reference count rather than as an
absolute score. Validation is the substantive difficulty, since
novelty has no ground truth; we approach it without one, through
axiomatic properties any such measure should satisfy, separation of
recognised landmark papers from matched controls, and correlation
with the disruption index and later citation impact. Our own
contribution is the model, the measurement of what a bibliography
does and does not determine, and the formulation and validation of
novelty as a conditional generative residual.

\end{document}
